\section{Discussion}
\label{sec:discussion}

% -----------------------------------------------------------------------
% Paragraph 1: Significance — H3S-class Tc at multi-anvil pressures
% -----------------------------------------------------------------------

The central finding of this work is that CsInH$_3$ achieves
$T_c \approx 200 \pm 50$~K at $P = 3$--$5$~GPa, matching the
experimental $T_c = 203$~K of H$_3$S~\cite{Drozdov2015} at a pressure
reduced by a factor of $30$. This pressure reduction transforms the
experimental accessibility of high-$T_c$ hydride superconductivity.
While 155~GPa requires diamond anvil cells with sample volumes of order
$10^{-4}$~mm$^3$, pressures of $3$--$5$~GPa are routinely achievable
with large-volume multi-anvil presses~\cite{walker1999}
that produce samples of order $1$~mm$^3$ --- a six-order-of-magnitude
increase in sample volume. This scale is sufficient for bulk transport
measurements, specific heat, and neutron scattering, all of which remain
impractical at megabar pressures. The quantum stabilization of CsInH$_3$
at 3~GPa (Sec.~\ref{sec:results}) provides a direct analogy to the
quantum stabilization of LaH$_{10}$ reported by
Errea~\textit{et~al.}~\cite{Errea2020}: in both cases, hydrogen
zero-point motion hardens a soft phonon mode that is unstable in the
harmonic approximation, enabling superconductivity at a pressure below
the classical stability threshold.

% -----------------------------------------------------------------------
% Paragraph 2: Lambda–omega_log tradeoff and Tc ceiling
% -----------------------------------------------------------------------

The $T_c(P)$ dome of CsInH$_3$ (Fig.~\ref{fig:tc_pressure}) reflects a
fundamental tradeoff between the electron-phonon coupling strength
$\lambda$ and the logarithmic average phonon frequency
$\omega_{\mathrm{log}}$. As pressure decreases from 10 to 3~GPa,
$\lambda_{\mathrm{anh}}$ grows from $\sim 1.9$ to $2.26$ due to phonon
softening, but the simultaneous decrease in $\omega_{\mathrm{log}}$ from
976~K to 831~K partially offsets the $T_c$ gain. This
$\lambda$--$\omega_{\mathrm{log}}$ tradeoff bounds the MXH$_3$
perovskite $T_c$ ceiling at $\sim 214$--$234$~K ($\mu^* = 0.13$--$0.10$)
within the Migdal-Eliashberg framework. The bound is consistent with the
analysis of Gao~\textit{et~al.}~\cite{gao2025ceiling}, who
showed that the product $\lambda \cdot \omega_{\mathrm{log}}$ is
constrained by lattice stability requirements at ambient pressure. At
moderate pressures ($3$--$10$~GPa), the constraint is relaxed --- softer
modes can persist without triggering structural collapse --- but not
sufficiently to reach $T_c = 300$~K. Achieving room-temperature
superconductivity in this chemical family would require
$\lambda_{\mathrm{anh}} > 3$ at $\omega_{\mathrm{log}} > 1000$~K, a
combination that anharmonic corrections preclude: SSCHA invariably
reduces $\lambda$ by $30$--$36\%$ in this class of materials.

% -----------------------------------------------------------------------
% Paragraph 3: Comparison with Du et al. and methodological caveats
% -----------------------------------------------------------------------

Our harmonic $T_c$ for CsInH$_3$ at 10~GPa ($245$~K at
$\mu^* = 0.13$) exceeds the value reported by
Du~\textit{et~al.}~\cite{Du2024} ($\sim 153$~K at
9~GPa) by approximately $60\%$, despite good agreement in $\lambda$
($2.35$ versus $\sim 2.4$). We trace this discrepancy to the spectral
function shape: our synthetic $\alpha^2 F(\omega)$ yields
$\omega_{\mathrm{log}} \approx 101$~meV, roughly $40\%$ higher than the
$\sim 65$~meV implied by the Du~\textit{et~al.} parameters obtained
with PBE~+~PAW and full DFPT. A $10$--$20\%$ difference from the
functional choice (PBEsol versus PBE) is expected for hydrogen phonon
frequencies~\cite{pickard2020}, but the
remaining discrepancy likely reflects the limitations of our synthetic
spectral function, which is calibrated against H$_3$S and LaH$_{10}$
rather than computed directly for CsInH$_3$ via density-functional
perturbation theory and Wannier interpolation.  If the lower
$\omega_{\mathrm{log}}$ of Du~\textit{et~al.} is confirmed by real EPW
calculations, the SSCHA-corrected $T_c$ would decrease to roughly
$140$--$180$~K --- still remarkably high for a $3$--$5$~GPa material,
but no longer matching H$_3$S. Resolving this $\omega_{\mathrm{log}}$
discrepancy is the single most important validation step for the
present predictions.

% -----------------------------------------------------------------------
% Paragraph 4: Limitations, outlook, and broader implications
% -----------------------------------------------------------------------

Several limitations of this work should be noted. First, all
$\alpha^2 F(\omega)$ spectra are synthetic (literature-calibrated);
absolute $T_c$ values carry a systematic uncertainty of $\sim 20$--$50\%$
that can only be resolved by full DFPT~+~EPW
calculations on the CsInH$_3$ structure. Second, the SSCHA treatment
employs eigenvector rotation calibrated against
H$_3$S~\cite{Errea2015} and YH$_6$~\cite{Belli2025} rather than a
self-consistent anharmonic electron-phonon calculation; the
$\lambda$ ratios are robust at the $\sim 5\%$ level, but the absolute
phonon frequencies inherit the synthetic baseline uncertainty. Third,
our isotropic Eliashberg treatment neglects anisotropic gap effects
that could modify $T_c$ by $5$--$15\%$ in the multi-sheeted perovskite
Fermi surface.
%
Looking forward, anisotropic Eliashberg calculations, together with an
assessment of whether CsInH$_3$ synthesized at 5~GPa can be quenched
to ambient pressure (kinetic metastability), would further establish the
viability of this material. Beyond CsInH$_3$, the chemical
pre-compression strategy motivates a systematic exploration of other
ternary hydride families --- Laves-phase, Heusler, and
clathrate-derived structures --- that may circumvent the
$\lambda$--$\omega_{\mathrm{log}}$ ceiling of the MXH$_3$ perovskites
and extend the reach of near-ambient-pressure high-$T_c$
superconductivity.

%% CITATIONS NEEDED (for gpd-bibliographer):
%% - walker1999 — Walker et al., multi-anvil press techniques
%% - gao2025ceiling — Gao et al. (2025), Tc ceiling analysis for hydrides
%% - Du2024 — Du et al., Advanced Science 11, 2408370 (2024), MXH3 perovskites
%% - pickard2020 — Review of functional dependence of hydride phonon frequencies
%% - Errea2020 — Errea et al., quantum stabilization of LaH10, Nature 2020
%% - Errea2015 — Errea et al., anharmonic effects in H3S, PRL 114, 157004 (2015)
%% - Belli2025 — Belli et al., anharmonic corrections in YH6, arXiv:2507.03383 (2025)
%% - Drozdov2015 — Drozdov et al., H3S 203K, Nature 525, 73 (2015)
